\documentclass[../DoAn.tex]{subfiles}
\begin{document}
\section{Kết luận}
Đồ án đã hoàn thành mục tiêu đề ra ở Chương~1: thiết kế và xây dựng thành công một hệ thống giám sát hành vi và phát hiện té ngã khép kín cho người cao tuổi, vận hành xuyên suốt từ thiết bị đeo nhúng đến ứng dụng web quản trị tập trung. Về phần cứng, thiết bị trên nền vi điều khiển ESP32-S3 thực hiện lấy mẫu cảm biến quán tính ổn định, nội suy mô hình học sâu ngay tại biên và truyền cảnh báo qua mạng 4G LTE bằng giao thức MQTT mà không phụ thuộc vào điện thoại trung gian. Về phần mềm, hệ thống máy chủ FastAPI với kiến trúc cơ sở dữ liệu kép (PostgreSQL cho siêu dữ liệu và InfluxDB cho dữ liệu chuỗi thời gian) cùng giao diện web thời gian thực cho phép nhân viên y tế giám sát đồng thời nhiều thiết bị và tiếp nhận cảnh báo té ngã với độ trễ đầu--cuối dưới một giây.

Quan trọng hơn, đồ án đã trả lời được câu hỏi nghiên cứu đặt ra ở Chương~1: một mạng học sâu hoàn toàn có thể đạt độ nhạy phát hiện té ngã cao mà vẫn suy luận trọn vẹn trên vi điều khiển tài nguyên hạn chế. Mô hình đề xuất cuối cùng --- \texttt{v30\_optimize}, một mạng tích chập tách kênh 1 chiều (Depthwise Separable 1D CNN) được đồng thiết kế theo tập lệnh tăng tốc ESP-NN --- đạt độ chính xác 95,30\% và Fall recall 97,50\% đo trực tiếp trên chip, với thời gian suy luận chỉ 56,7~ms (khoảng 11\% chu kỳ cửa sổ trượt 500~ms) và dung lượng xấp xỉ 55,8~KB. Ba đóng góp kỹ thuật nổi bật được trình bày ở Chương~5 gồm: (i)~chiến lược thiết kế nhãn 5 lớp kết hợp gán nhãn theo sự kiện nhằm giảm thiểu cảnh báo giả; (ii)~quá trình đồng thiết kế kiến trúc hướng phần cứng, chắt lọc nguyên lý mô hình hóa thời gian của mạng TCN vào một CNN thân thiện ESP-NN; và (iii)~cơ chế lấy mẫu IMU chính xác bằng ngắt phần cứng PCNT.

So với các hướng tiếp cận hiện có đã khảo sát ở Chương~2, đồ án lấp được khoảng trống công nghệ giữa hai thái cực: các mô hình học sâu chính xác cao nhưng buộc phải chạy trên Cloud/PC, và các thiết bị đeo tự tính toán nhưng bị giới hạn ở các thuật toán ngưỡng hoặc máy học cơ sở dễ phát sinh báo động giả. Bằng việc đưa độ chính xác của học sâu xuống chạy trực tiếp trên vi điều khiển giá thành thấp, hệ thống vừa bảo toàn quyền riêng tư (không sử dụng camera, không truyền dữ liệu thô lên máy chủ), vừa hoạt động độc lập và liên tục trên mạng di động.

Bên cạnh các kết quả đạt được, đồ án vẫn còn một số hạn chế. Mô hình được đánh giá chủ yếu trên bộ dữ liệu công khai SisFall kết hợp một tập dữ liệu tự thu quy mô nhỏ, chưa được kiểm chứng trên số lượng lớn người cao tuổi trong điều kiện thực tế. Phiên bản \texttt{v30\_optimize} hiện là phương án đề xuất; việc thay thế hẳn mô hình \texttt{v25} đang vận hành trong firmware sản phẩm và hoàn thiện cơ chế ngủ tiết kiệm điện vẫn là bước triển khai kế tiếp. Bài học lớn nhất rút ra là nguyên tắc \textit{thiết kế thuận theo phần cứng} (hardware-aware design): một mô hình chính xác hơn trên máy chủ chưa chắc đã khả dụng trên thiết bị biên --- ràng buộc của tập lệnh tăng tốc phải được đưa vào ngay từ khâu thiết kế kiến trúc, thay vì tối ưu hóa sau khi đã huấn luyện.

\section{Hướng phát triển}
Để hoàn thiện hệ thống thành một sản phẩm hoàn chỉnh, trước mắt cần tích hợp chính thức mô hình \texttt{v30\_optimize} vào firmware sản phẩm, hoàn thiện cơ chế lấy mẫu PCNT kèm chế độ ngủ tiết kiệm điện (Automatic Light Sleep) đã đặc tả ở Mục~5.3, đồng thời mở rộng tập dữ liệu huấn luyện bằng dữ liệu vận động thực tế của người cao tuổi thu thập qua chính pipeline thu thập dữ liệu của hệ thống. Bên cạnh các công việc hoàn thiện này, đồ án xác định một số hướng nâng cấp giàu tiềm năng sau đây.

\textbf{Tối ưu năng lượng ở khâu cảm biến bằng IMU có FIFO và ngắt thông minh.} Cơ chế lấy mẫu PCNT hiện tại tuy bảo đảm tần số chính xác nhưng phụ thuộc vào hoạt động liên tục của ngoại vi, gây khó khăn khi đưa CPU vào chế độ Light Sleep --- rào cản trực tiếp cho việc tiết kiệm điện sâu. Hướng khắc phục là chuyển sang các IMU thế hệ mới (như Bosch BMI270, STMicroelectronics LSM6DSO hay TDK InvenSense ICM-42688) tích hợp sẵn bộ đệm FIFO phần cứng cùng ngắt theo ngưỡng đầy (FIFO watermark interrupt): cảm biến tự tích lũy mẫu trong FIFO và chỉ đánh thức vi điều khiển khi bộ đệm đầy tới mức định trước, cho phép CPU ngủ giữa các đợt thay vì phải thức liên tục ở tần số 100~Hz. Sâu hơn nữa, các IMU này còn hỗ trợ ngắt phát hiện chuyển động ở mức phần cứng (wake-on-motion, any-motion, significant-motion): thiết bị có thể chuyển vào chế độ Deep Sleep khi người dùng nằm hoặc ngồi yên trong thời gian dài và chỉ ``bừng tỉnh'' khi phát sinh chuyển động đáng kể. Đặc thù người cao tuổi phần lớn thời gian ở trạng thái tĩnh khiến cơ chế này có tiềm năng kéo dài tuổi thọ pin lên nhiều lần.

\textbf{Tối ưu năng lượng ở khâu truyền thông di động.} Module A7680C hiện sử dụng chuẩn LTE Cat-1 duy trì kết nối vô tuyến (trạng thái RRC Connected) trong khoảng 10 giây sau mỗi lần publish bản tin --- đây là thành phần tiêu thụ dòng điện chủ yếu đối với một thiết bị chỉ gửi tin thưa thớt. Hướng tối ưu là chuyển sang các module LTE-M (Cat-M1) hoặc NB-IoT vốn được thiết kế chuyên cho IoT công suất thấp và hỗ trợ tín hiệu RAI (Release Assistance Indication): ngay sau khi gửi xong gói tin cảnh báo, thiết bị báo cho mạng biết không còn dữ liệu để giải phóng kết nối tức thì, loại bỏ hoàn toàn khoảng ``treo'' 10 giây tốn điện của chuẩn Cat-1.

\textbf{Đa dạng hóa giao thức cho môi trường viện dưỡng lão cố định.} Trong môi trường triển khai cố định như viện dưỡng lão, có thể bổ sung các chuẩn truyền thông mạng lưới (mesh) công suất thấp như Zigbee hoặc Thread/Matter, kết hợp cùng một cổng kết nối (gateway) trung tâm. So với phương án mỗi thiết bị mang một module di động riêng, giải pháp mesh cục bộ giảm đáng kể chi phí và điện năng tiêu thụ trên mỗi nút mạng, đồng thời chuẩn Matter/Thread mở ra khả năng tương tác liền mạch với hạ tầng nhà thông minh và các thiết bị y tế khác trong cùng cơ sở.

\textbf{Tích hợp cảm biến sinh hiệu và thiết kế PCB hướng sản phẩm.} Để tiến gần tới một sản phẩm thương mại hoàn chỉnh, một hướng phát triển quan trọng là thiết kế bo mạch in (PCB) tối ưu, tích hợp thêm module đo điện tâm đồ (ECG) hoặc nhịp tim quang học (PPG). Điều này cho phép hệ thống phát hiện sớm các biến cố tim mạch như nhồi máu cơ tim hay đột quỵ --- vốn thường chính là nguyên nhân gốc rễ dẫn tới té ngã ở người cao tuổi. Khi nắm bắt được cả nguyên nhân (sự cố tim mạch) lẫn hậu quả (cú ngã), hệ thống không những cảnh báo kịp thời hơn mà còn cung cấp một bức tranh sức khỏe toàn diện, nâng tầm thiết bị từ một cảm biến phát hiện ngã đơn thuần thành một thiết bị đeo theo dõi sức khỏe đa năng.
\end{document}
